\begin{lemma}[Restricting to the full domain is the identity]
  Let $E$ be a metric space and $A \in \Theta(E)$ (\noderef{Curve}). Then
  \[
    A|_{[0,\length(A)]} = A
    ,
  \]
  i.e.\ $\mathrm{curveRestrict}(A,0,\length(A)) = A$ (\noderef{curveRestrict}).

  \paragraph{Why this is the enabler for the arc repair.} Every one of
  \noderef{RestrictLSI}, \noderef{RestrictDeltaEpsNAtBreakpoints}, \noderef{RestrictPiecewiseGeodesic}
  and \noderef{RestrictMeasureFormula} is stated about a sub-window $\mathrm{curveRestrict}(\gamma,u,u')$
  of a \emph{generic} ambient curve $\gamma$. Taking the ambient curve to be $A :=
  \mathrm{curveArc}(\gamma,t,t')$ (\noderef{curveArc}) itself, this lemma identifies
  $\mathrm{curveRestrict}(A,0,\length(A))$ with $A$, so any of those four lemmas applied at
  $t:=0,\ t':=\length(A)$ to ambient $A$ gives a statement directly about $A$'s own l.s.i., $\delta
  \epsilon n$-count, piecewise-geodesic, or measure property, with no need to restate any of the four
  for the wrap case: $A$ is an ordinary $\Theta(E)$-curve whether or not the arc it names wraps
  around the closed curve $\gamma$.
\end{lemma}
\begin{proof}
  Write $B := \mathrm{curveRestrict}(A,0,\length(A))$. By \noderef{curveRestrict}'s defining
  formula (with $a:=0$, $b:=\length(A)$), $\length(B) = \length(A) - 0 = \length(A)$, and for every
  $u \in \mathbb{R}$,
  \[
    B(u) = A\bigl(0 + \min(\length(A)-0,\ \max(0,u))\bigr) = A\bigl(\min(\length(A),\max(0,u))\bigr)
    .
  \]
  By $A$'s own clamping property (\noderef{Curve}), $A(u) = A(\min(\length(A),\max(0,u)))$ for
  every $u$. Hence $B(u) = A(u)$ for every $u \in \mathbb{R}$, i.e.\ the underlying maps of $B$ and
  $A$ agree identically.

  So $B$ and $A$ are two elements of $\Theta(E)$ with the same length ($\length(B) = \length(A)$)
  and the same underlying map ($B(u) = A(u)$ for all $u$). The remaining structural data attached
  to a curve in $\Theta(E)$ (\noderef{Curve}) -- non-negativity of the length, the clamping
  condition, and the unit-speed condition -- are each a proof of a proposition depending only on
  the length and the underlying map, hence, once those two agree, the corresponding proof
  components of $B$ and $A$ are proofs of the same proposition and $B$ and $A$ coincide as elements
  of $\Theta(E)$: $B = A$.
\end{proof}
